% !TEX root = ../notes.tex

\documentclass[../notes.tex]{subfiles}

\begin{document}

\section{March 18}
Here we go.

\subsection{The Pauli Exclusion Principle}
Here is our definition.
\begin{definition}[Pauli exclusion]
	If the space of states of a single electron is $V$, then the space of sates of $n$ electrons is $\land^nV$.
\end{definition}
Notably, if states were actual vectors in classical mechanics, then we would expect that the collection of spaces of two vectors would be $\op{Sym}^2V$ because symmetry should be preserved if the electrons are indistinguishable. However, in quantum mechanics, our states only matter up to phase, so it would be okay for permutations to act up to a scalar. There are two such one-dimensional characters of the symmetric group, which correspond to bosons and fermions.
\begin{definition}[boson]
	A \textit{boson} is a kind of state in a space $V$ for which the space of $n$ bosons is $\op{Sym}^nV$.
\end{definition}
\begin{definition}[fermion]
	A \textit{fermion} is a kind of state in a space $V$ for which the space of $n$ bosons is $\land^nV$.
\end{definition}
Thus, the Pauli exclusion principle is just the physical observation that electrons and fermions.
\begin{example}
	The space $V_n$ of electrons at energy level $-\frac1{2n^2}$ has dimension $2n^2$. Thus, the space of $k$ electrons is $\land^kV_n$; notably, there cannot be more than $2n^2$ electrons at a given level. For example, at the first energy level $-1/2$, we are allowed at most two energy levels, which correspond to the two electrons in helium. The next energy level has eight permissible electrons.
\end{example}
Frequently, the presence of repeated eigenvalues of some matrix is due to the presence of an underlying representation. For example, the adjacency matrix of a graph $G$ will frequently have repeated eigenvalues when $G$ has a large automorphism group. There is a similar notion of ``degeneracy'' in physics due to our repeated energy levels. However, note that
\[W_n=2L_1\oplus2L_3\cdots\oplus2L_{2n-3}\oplus L_{2n-1}\]
as a representation of $\mf{sl}_2\oplus\mf{sl}_2$. The explanation for the many $2$s is due to the presence of an additional symmetry of the hydrogen atom: setting $\vec r\coloneqq(x,y,z)$ and $\vec p\coloneqq(-i\del_x,-i\del_y,-i\del_z)$, there is an angular momentum operator $\vec L\coloneqq\vec r\times\vec p$. The components of $\vec L$ define the action of $\mf{so}(3)$, so because ${\vec p}^2=\vec p\cdot\vec p$ is the negative Laplacian, which is $\mf{so}(3)$-invariant, we see $[\vec L,{\vec p}^2]=0$.

We now define the second-order differential operator
\[\vec A_0\coloneqq\frac12\left(\vec p\times\vec L-\vec L\times\vec p\right),\]
which automatically commutes with $\vec p^2$. We would like for $\vec A_0$ to also commute with the Hamiltonian $H=-\frac12\Delta+U(r)=-\frac12p^2+U(r)$. To this end, we modify $\vec A_0$ to be $\vec A\coloneqq\vec A_0+\varphi(r)\vec r$, and we would like for $[\vec A,H]=0$. On the homework, we will show that such a $\varphi$ exists if and only if $U(r)=\frac Cr+D$ (which is the Coulomb potential), and then one can compute $\varphi$. The moral is that the components of $\vec A$ define a third $\mf{sl}_2$-triple acting on $L^2\left(\RR^3\right)$ and commuting with our other two $\mf{sl}_2=\mf{so}(3)$ acting on $W_n$.
\begin{remark}
	We call this extra $\mf{sl}_2$ a ``hidden symmetry.'' Such things also exist classically: for example, Kepler's laws are periodic only due to the special nature of the potential. Small adjustments to the potential no longer have 
\end{remark}
It now turns out that $W_n=L_{n-1}\boxtimes L_{n-1}$ as a representation of the combined $\mf{sl}_2\oplus\mf{sl}_2$; notably, it is irreducible, and one finds the decomposition $L_0\oplus L_2\oplus\cdots\oplus L_{2n}$ by decomposing the box product against the diagonal $\mf{sl}_2\subseteq\mf{sl}_2\oplus\mf{sl}_2$. Accounting for spin, we get $L_{n-1}\boxtimes L_{n-1}\boxtimes L_1$.

\subsection{The Automorphism Group}
We have some time before spring break, so let's say something about our next topic. Let's start by recalling the following facts.
\begin{proposition}
	Fix a complex semisimple Lie algebra $\mf g$. Then $\op{Aut}(\mf g)$ is a complex Lie group with Lie algebra $\mf g$.
\end{proposition}
\begin{proof}
	Note that $\op{Aut}(\mf g)$ is a closed Lie subgroup of $\op{GL}(\mf g)$. One can check that the Lie algebra of $\op{Aut}(\mf g)$ is given by the derivations $\mf g\to\mf g$. However, a Lie algebra gives a large source of derivations: there is a map $\mf g\to\op{Der}(\mf g)$ given by sending $X$ to the adjoint action. For a semisimple algebra, this map is an isomorphism of Lie algebras, so the result follows.
\end{proof}
\begin{definition}[adjoint group]
	Fix a complex semisimple Lie algebra $\mf g$. Then we define the \textit{adjoint group} $G_{\mathrm{ad}}$ to be $\op{Aut}(\mf g)^\circ$.
\end{definition}
\begin{example}
	For $\mf g=\mf{sl}_n$, one can check that $G_{\mathrm{ad}}$ is $\op{PGL}_n(\CC)$.
\end{example}
We also had the following classification of Cartan subalgebras.
\begin{proposition}
	Fix a complex semisimple Lie algebra $\mf g$. Then all Cartan subalgebras $\mf h\subseteq\mf g$ are conjugate in the sense that $\op{Aut}(\mf g)^\circ$ acts transitively on $\mf g$.
\end{proposition}
\begin{proof}
	Omitted.
\end{proof}
We are going to get quite a bit of mileage out of a choice of Cartan $\mf h\subseteq\mf g$ and lift to $G_{\mathrm{ad}}$.
\begin{lemma}
	Fix a Cartan subalgebra $\mf h$ of a complex semisimple Lie algebra $\mf g$. Then $\exp(\mf h)\subseteq G_{\mathrm{ad}}$ is a Lie subgroup isomorphic to $(\CC^\times)^{\op{rank}\mf g}$.
\end{lemma}
\begin{proof}
	Let $H\subseteq G$ be the Lie subgroup corresponding to $\mf h$. Because $\mf h$ is commutative, we see that $H$ is commutative, so $\exp\colon\mf h\to H$ is a group homomorphism. Because $\exp$ is also a local diffeomorphism, we conclude that its image contains an open neighborhood of $H$, so $\exp$ is surjective (because $H$ is connected).

	Now, for $X\in\mf h$, we see that $\op{ad}_X$ acts by $0$ on $\mf h$ and by $\alpha(X)$ on the root spaces $\mf g_\alpha$. This also describes the differential action of $H\subseteq G_{\mathrm{ad}}$ on $\mf g$, so we may take the exponential in $G_{\mathrm{ad}}\subseteq\op{GL}(\mf g)$ to see that $\exp(X)$ must act by $1$ on $\mf h$ and by $e^{\alpha(X)}$ on $\mf g_\alpha$. In particular, the action of $H$ on $\mf g$ diagonalizes (via the root space decomposition of $\mf g$).
	
	We also see that the kernel of $\exp\colon\mf h\to H$ is given by those $X\in\mf h$ for which $\alpha(X)\in2\pi i\ZZ$ for all roots $\alpha$, so it follows that $H$ is isomorphic to $\mf h/2\pi iP^\lor$. Taking a quotient of a complex vector space by a lattice produces a power of $\CC^\times$, so we are done. (Indeed, we merely have to take powers of the isomorphism $\exp\colon\CC/2\pi i\ZZ\to\CC^\times$.)
\end{proof}
Let's use this $H$ for fun and profit.
\begin{proposition} \label{prop:compute-weyl-group}
	Fix a Cartan subalgebra $\mf h$ of a complex semisimple Lie algebra $\mf g$, and set $H\coloneqq\exp(\mf h)$ in $G_{\mathrm{ad}}$. Then the normalizer of $N(H)$ of $H$ in $G_{\mathrm{ad}}$ coincides with the stabilizer of $\mf h$ and contains $H$ as a normal subgroup. In fact, $N(H)/H$ is isomorphic to the Weyl group.
\end{proposition}
\begin{proof}
	Set $G\coloneqq G_{\mathrm{ad}}$ for brevity. Certainly $H$ is normal in $N(H)$ by construction of $N(H)$. Furthermore, recall that $G_{\mathrm{ad}}$ acts transitively on the collection of Cartan subalgebras, and $g\in G_{\mathrm{ad}}$ fixes $\mf h$ (by adjoint action) if and only if $g$ fixes $H$ (also by adjoint action!) by the uniqueness of the connected Lie subgroups attached to Lie subalgebras. It follows that $N(H)$ is the stabilizer of $\mf h$.
	
	Thus, the main content of the proof lies in the last statement. We will now only show the last statement. The general approach will be to first lift the Weyl group to $G$, and then we will show that $H$ and the lifted Weyl group generate $N(H)$. Let's start with the lifting. For each $i$, there is a map $\eta_i\colon\mf{sl}_2\to\mf g$ given by sending $e\mapsto e_i$ and $f\mapsto f_i$ and $h\mapsto h_i$. Because $\op{SL}_2(\CC)$ is simply connected, this lifts (uniquely) to a map $\eta_i\colon\op{SL}_2(\CC)\to G_{\mathrm{ad}}$. This $\eta_i$ let us define
	\[w_i\coloneqq\eta_i\left(\begin{bmatrix}
		0 & 1 \\ -1 & 0
	\end{bmatrix}\right).\]
	We will use this element as a lift of the simple reflection $s_i\in W$. We can now lift the rest of the Weyl group. Given $w\in W$, we may pick a decomposition $w=s_{i_1}\cdots s_{i_n}$ into simple reflections, and we can define $\widetilde w\coloneqq w_{i_1}\cdots w_{i_n}$.\footnote{It turns out that $\widetilde w$ does not depend on the choice of decomposition of $w$, but we will not need this.} By construction, the adjoint action of $w_i$ on $\mf h$ by $s_i$: indeed, it acts by reflecting $\alpha_i$'s root space.
	
	We emphasize that our lifting is not a group homomorphism: for example, $w_i^2=\eta_i(-1)$, which may or may not be the identity. On the other hand, note that the differential action of $w_i$ fixes $\mf h$, so it lives in $N(H)$. However, it will turn out to be a homomorphism ``up to $H$.'' Let's make this precise. For any $w,u\in W$, we do know that $\widetilde{wu}=\widetilde w\widetilde uh$ for some unique $h\in G_{\mathrm{ad}}$; we claim that actually $h\in H$. Well, $\widetilde{wu}$ and $\widetilde w\widetilde u$ both preserve $\mf h$ and even have the same action on the roots, so it follows that $h$ preserves all root spaces. Accordingly, we are granted complex numbers $b_j$ for which $h|_{\mf g_{\alpha_j}}=e^{b_j}$. By comparing the action on $\mf g$, we conclude that
	\[h=e^{\sum_jb_j\omega_j^\lor},\]
	which one shows by induction on the height of the root spaces. (Note that the action of $h$ on $\mf g$ is enough to know $h$ because $G_{\mathrm{ad}}\subseteq\op{GL}(\mf g)$.) Thus, $h\in H$.
	
	It follows that the elements $\widetilde w_i$ together with $H$ generate a subgroup $N$ of $G$ sitting in a short exact sequence
	\[1\to H\to N\to W\to1.\]
	It remains to show that $N=N(H)$, which will complete the proof. Well, choose $g\in N(H)$, which we would like to move to $N$. But then $\{g(\alpha_i)\}_i$ is a collection of positive roots, but all systems of positive roots are conjugate by the Weyl group, so we can adjust $g$ by an element $\widetilde w\in N$ so that $g$ acts by a permutation of the simple roots.
	
	We are now in the situation where $g$ acts by an automorphism on the Dynkin diagram. If we could show that this automorphism is trivial, then it would follow that $g$ merely scales all the root spaces (with no permutation), so we are able to conclude that $g\in H$ using the argument above. However, we claim that all of $G_{\mathrm{ad}}$ acts trivially on the Dynkin diagram (because the automorphism group of the Dynkin diagram is discrete). It is enough to check that $G_{\mathrm{ad}}$ acts trivially on the fundamental representations of $\mf g$: indeed, $G_{\mathrm{ad}}$ permutes the Verma modules according to its action on $\mf h$, which then has quotient to the needed fundamental representations. But $G_{\mathrm{ad}}$ needs to act continuously on the characters, which distinguish the fundamental representations, so we conclude that $G_{\mathrm{ad}}$ acts trivially on the characters and thus the fundamental representations.
	% \todo{}
	% \footnote{Professor Etingof made some argument about elements of the group acting on representations, and $G_{\mathrm{ad}}$ needs to act trivially on the isomorphism classes of representations by looking at the characters. Alternatively, one can use the fact that representations of $\mf g$ lift uniquely to the universal cover of $G_{\mathrm{ad}}$, so conjugation is not allowed to do anything.} Thus, the permutation is trivial, so $\widetilde w^{-1}g$ preserves $\mf h$ and all the root spaces, so it is in $H$ as before, so $g$ lives in $N$.
\end{proof}
\begin{remark}
	The extension
	\[1\to H\to N(H)\to W\to1\]
	does not always split. Indeed, this does not split for $\mf g=\mf{sl}_2(\CC)$: one sees that $N(H)$ is generated by the central torus of diagonal elements $\op{diag}(t,1/t)$ and the element $w=\begin{bsmallmatrix}
		0 & 1 \\ -1 & 0
	\end{bsmallmatrix}$, which satisfies $wtw^{-1}=t^{-1}$. Notably, $(wt)^2=-1$ for any $t$, but this is not possible for elements in $W\rtimes H$.
\end{remark}
We can upgrade this discussion to compute $\op{Aut}(\mf g)$ in full.
\begin{proposition} \label{prop:aut-lie-alg}
	Let $\mf g$ be a complex semisimple Lie algebra with Dynkin diagram $D$. Then there is a split short exact sequence
	\[1\to G_{\mathrm{ad}}\to\op{Aut}(\mf g)\to\op{Aut}(D)\to1.\]
\end{proposition}
\begin{proof}
	To begin, note that the automorphisms of the Dynkin diagram embed into $\op{Aut}(\mf g)$ by using the Serre presentation of $\mf g$. Accordingly, there is a map
	\[\xi\colon\op{Aut}(D)\ltimes G_{\mathrm{ad}}\to\op{Aut}(\mf g).\]
	Indeed, certainly $G_{\mathrm{ad}}$ is normal in $\op{Aut}(\mf g)$ because it is the connected component. Also, we may quickly check that $\xi$ is injective: certainly $\xi$ is injective on $\op{Aut}(D)$. Thus, if we have a nontrivial element $x\in\ker\xi$, the image of $x$ in $\op{Aut}(D)$ is not the identity. But then $x$ permutes the fundamental representations nontrivially, which does not happen for elements of $G_{\mathrm{ad}}$.

	It remains to show that $\xi$ is surjective. Well, choose $g\in\op{Aut}(\mf g)$. Because any two Cartan subalgebras are conjugate by $G_{\mathrm{ad}}$, we may assume that $g$ fixes $\mf h$. We now redo our previous argument: the collection $\{g(\alpha_i)\}$ is a new collection of simple roots, so there is an automorphism $\sigma$ of the Dynkin diagram with $g(\alpha_i)=\alpha_{\sigma(i)}$. Adjusting $g$ by $\sigma^{-1}$, we may further assume that $g$ fixes the simple roots as well. We conclude that $g\in H$ using the argument from before.
\end{proof}

\end{document}